\begin{codebox}
\codeline{\textcolor{cmtgray}{\#\ 哲学本体论与工程本体论的对照}}
\codeline{\textcolor{cmtgray}{\#\ From\ Philosophical\ Ontology\ to\ Engineering\ Ontology}}
\codeline{\textcolor{cmtgray}{\#}}
\codeline{\textcolor{cmtgray}{\#\ 本体论（Ontology）一词源于哲学，AI与知识工程借用并重新定义了它}}
\codeblank{}
\codeline{\textcolor{cmtgray}{\#\ ==============================}}
\codeline{\textcolor{cmtgray}{\#\ 1.\ 哲学本体论的追问}}
\codeline{\textcolor{cmtgray}{\#\ ==============================}}
\codeline{\textcolor{cmtgray}{\#\ 哲学本体论（首字母大写的\ Ontology）研究"存在之为存在"：}}
\codeline{\textcolor{cmtgray}{\#\ \ \ 世界的本源是什么？}}
\codeline{\textcolor{cmtgray}{\#\ \ \ 一个事物凭什么是"它自己"？}}
\codeline{\textcolor{cmtgray}{\#\ \ \ 共相（普遍概念）是否真实存在？}}
\codeblank{}
\codeline{\textcolor{cmtgray}{\#\ 亚里士多德的范畴论给出了最早的"顶层分类"：}}
\codeline{\textcolor{cmtgray}{\#\ \ \ 实体（Substance）、数量（Quantity）、性质（Quality）、}}
\codeline{\textcolor{cmtgray}{\#\ \ \ 关系（Relation）、地点（Place）、时间（Time）……}}
\codeblank{}
\codeline{\textcolor{cmtgray}{\#\ 特征：终极追问无法收敛，知识永远处于开放的不确定性中}}
\codeline{\textcolor{cmtgray}{\#\ ——\ 这类问题"无法落地、不可计算"，不是工程本体论的研究对象}}
\codeblank{}
\codeline{\textcolor{cmtgray}{\#\ ==============================}}
\codeline{\textcolor{cmtgray}{\#\ 2.\ 工程本体论的转向}}
\codeline{\textcolor{cmtgray}{\#\ ==============================}}
\codeline{\textcolor{cmtgray}{\#\ 知识工程领域（小写复数的\ ontologies）完成了一次实用主义转向：}}
\codeline{\textcolor{cmtgray}{\#\ 不再追问存在的本质，而是约定——}}
\codeblank{}
\codeline{\textcolor{cmtgray}{\#\ \ \ 在一个领域内，存在哪些概念（类）？}}
\codeline{\textcolor{cmtgray}{\#\ \ \ 概念之间有什么关系（属性）？}}
\codeline{\textcolor{cmtgray}{\#\ \ \ 什么样的陈述是合法的（公理与约束）？}}
\codeblank{}
\codeline{\textcolor{cmtgray}{\#\ Gruber\ (1993)\ 的经典定义：}}
\codeline{\textcolor{cmtgray}{\#\ \ \ "An\ ontology\ is\ an\ explicit\ specification\ of\ a\ conceptualization."}}
\codeline{\textcolor{cmtgray}{\#\ \ \ 本体是概念化的显式规范}}
\codeblank{}
\codeline{\textcolor{cmtgray}{\#\ Studer\ et\ al.\ (1998)\ 的综合定义包含四个关键词：}}
\codeline{\textcolor{cmtgray}{\#\ \ \ 概念化（Conceptualization）——\ 对世界某侧面的抽象模型}}
\codeline{\textcolor{cmtgray}{\#\ \ \ 显式（Explicit）\ \ \ \ \ \ \ \ \ \ ——\ 概念与约束被明确定义，而非隐含}}
\codeline{\textcolor{cmtgray}{\#\ \ \ 形式化（Formal）\ \ \ \ \ \ \ \ \ \ ——\ 机器可读、可推理，而非自然语言描述}}
\codeline{\textcolor{cmtgray}{\#\ \ \ 共享（Shared）\ \ \ \ \ \ \ \ \ \ \ \ ——\ 团体共识，而非个人观点}}
\codeblank{}
\codeline{\textcolor{cmtgray}{\#\ ==============================}}
\codeline{\textcolor{cmtgray}{\#\ 3.\ 对照：同一个"设备"概念}}
\codeline{\textcolor{cmtgray}{\#\ ==============================}}
\codeline{\textcolor{cmtgray}{\#\ 哲学视角的追问：}}
\codeline{\textcolor{cmtgray}{\#\ \ \ "设备"作为人造物（Artifact）的本质是什么？}}
\codeline{\textcolor{cmtgray}{\#\ \ \ 一台车床更换全部零件后还是原来那台车床吗？（忒修斯之船）}}
\codeblank{}
\codeline{\textcolor{cmtgray}{\#\ 工程视角的约定：}}
\codeline{Equipment\ \(\sqsubseteq\)\ PhysicalObject\ \ \ \ \ \ \ \ \ \ \ \ \ \ \ \#\ 设备是物理对象}
\codeline{Equipment\ \(\sqsubseteq\)\ \(\exists\)hasSerialNumber.string\ \ \ \ \ \ \#\ 每台设备有序列号（同一性判据）}
\codeline{CNCLathe\ \(\sqsubseteq\)\ CNCMachine\ \(\sqsubseteq\)\ Equipment\ \ \ \ \ \ \ \ \#\ 分类层次}
\codeline{Equipment\ \(\sqcap\)\ Material\ \(\sqsubseteq\)\ \(\bot\)\ \ \ \ \ \ \ \ \ \ \ \ \ \ \ \ \ \#\ 设备与材料不相交}
\codeblank{}
\codeline{\textcolor{cmtgray}{\#\ 工程本体论不回答"本质"问题，}}
\codeline{\textcolor{cmtgray}{\#\ 而是给出可检验、可推理、可共享的形式化约定：}}
\codeline{\textcolor{cmtgray}{\#\ 序列号就是这个信息系统中设备的同一性判据——问题被"约定"解决}}
\codeblank{}
\codeline{\textcolor{cmtgray}{\#\ ==============================}}
\codeline{\textcolor{cmtgray}{\#\ 4.\ 知识熵减的视角}}
\codeline{\textcolor{cmtgray}{\#\ ==============================}}
\codeline{\textcolor{cmtgray}{\#\ 企业知识的自然状态是高熵的：}}
\codeline{\textcolor{cmtgray}{\#\ \ \ 分散在文档、数据库、员工头脑中}}
\codeline{\textcolor{cmtgray}{\#\ \ \ 同一概念多种叫法，同一名称多种含义}}
\codeline{\textcolor{cmtgray}{\#\ \ \ 规则隐含在经验里，无法被系统校验}}
\codeblank{}
\codeline{\textcolor{cmtgray}{\#\ 本体工程的本质是知识熵减：}}
\codeline{\textcolor{cmtgray}{\#\ \ \ 隐性\ \(\rightarrow\)\ 显式：把头脑中的概念写成形式化定义}}
\codeline{\textcolor{cmtgray}{\#\ \ \ 分散\ \(\rightarrow\)\ 统一：建立全企业共享的概念词汇表}}
\codeline{\textcolor{cmtgray}{\#\ \ \ 模糊\ \(\rightarrow\)\ 精确：用公理排除非法状态}}
\codeline{\textcolor{cmtgray}{\#\ \ \ 静态\ \(\rightarrow\)\ 可推理：让机器从已有知识推导新知识}}
\codeblank{}
\codeline{\textcolor{cmtgray}{\#\ 熵减的产出：}}
\codeline{\textcolor{cmtgray}{\#\ \ \ 形式化本体（TBox）+\ 实例数据（ABox）}}
\codeline{\textcolor{cmtgray}{\#\ \ \ =\ 可查询、可推理、可复用、可解释的知识结构}}
\codeblank{}
\codeline{\textcolor{cmtgray}{\#\ ==============================}}
\codeline{\textcolor{cmtgray}{\#\ 5.\ 工程判据}}
\codeline{\textcolor{cmtgray}{\#\ ==============================}}
\codeline{\textcolor{cmtgray}{\#\ 一个"好的"工程本体不以哲学正确性为判据，而以下列工程性质为判据：}}
\codeline{\textcolor{cmtgray}{\#\ \ \ 一致性（Consistency）\ \ ——\ 公理之间无矛盾，推理器可检验}}
\codeline{\textcolor{cmtgray}{\#\ \ \ 完备性（Completeness）\ ——\ 覆盖领域内需要回答的问题（见第3章能力问题）}}
\codeline{\textcolor{cmtgray}{\#\ \ \ 简约性（Minimality）\ \ \ ——\ 不引入与任务无关的概念}}
\codeline{\textcolor{cmtgray}{\#\ \ \ 可扩展性（Extensibility）——\ 新概念可平滑加入而不推翻已有结构}}
\codeline{\textcolor{cmtgray}{\#\ \ \ 共享性（Consensus）\ \ \ \ ——\ 领域专家认可并实际使用}}
\end{codebox}
